\section{Companion results on algebraic correspondences}\label{sec:companion}

We now separate three different assertions suggested by the finite-field
realization problem. Some varieties admit explicit retractions in
rational Chow motives. For other varieties, an abelian motive with the
right Frobenius realization can be constructed, but an algebraic
correspondence from the original variety remains to be found. Finally,
for a specified target, that remaining assertion is equivalent to a
precise mixed Tate-conjecture statement. None of these distinctions is
removed by the descent calculation above.

Write $\mathsf h(X)$ for the motive of a smooth projective variety and
$\mathbb L=\Q(-1)$ for the Lefschetz motive, with Frobenius eigenvalue
$q$ over $\F_q$. All correspondences have rational coefficients.
We use the convention
\[
 \operatorname{Hom}(\mathsf h(X),\mathsf h(Y))
       =\operatorname{CH}^{\dim X}(X\times Y)_{\Q}
\]
for untwisted motives; the induced map preserves cohomological degree.
The first two constructions below take place in Chow motives, so their
identities hold modulo rational equivalence. The subsequent target and
comparison statements take place modulo homological equivalence and
assert identities in every $\ell$-adic realization, $\ell\ne p$.
They do not assert Chow isomorphisms.

\subsection{Retractions and finite-field descent}

Two elementary operations turn geometric maps over an extension
field into motivic summands over the original field.

\begin{lemma}\label{lem:retractions}
Let $X$ and $V$ be smooth projective varieties of the same dimension.
If $\phi:V\to X$ is a proper generically finite morphism of degree $a$,
then $\mathsf h(X)$ is a Chow summand of $\mathsf h(V)$, with inclusion
$\phi^*$ and retraction $a^{-1}\phi_*$. If $Z\subset Y$ is smooth of
codimension two, then
\[
 \mathsf h(\operatorname{Bl}_Z Y)
       \simeq \mathsf h(Y)\oplus\mathsf h(Z)\otimes\mathbb L.
\]
\end{lemma}

\begin{proof}
The projection formula gives $\phi_*\phi^*=a$ on Chow groups after
product with every smooth projective variety. Equivalently, composition
of the two graph correspondences is $a\Delta_X$, which proves the
first assertion. For the second, let $b:\widetilde Y\to Y$ be the
blowup and $j:E\hookrightarrow\widetilde Y$, $\rho:E\to Z$ its
exceptional divisor and projection. The maps are
\[
 (b^*,j_*\rho^*),\qquad (b_*,-\rho_*j^*).
\]
The identities $b_*b^*=1$, $b_*j_*\rho^*=0$,
$\rho_*j^*b^*=0$, and
$-\rho_*j^*j_*\rho^*=1$ follow from the projection formula and
$\mathcal O_E(E)=\mathcal O_E(-1)$. The blowup formula says that
these maps exhaust the Chow groups, after every product; hence they
are inverse motivic maps. See \citet[Theorem 2.8(3), p.~1472,
and Corollary 3.10, p.~1484]{Jiang2023}.
\end{proof}

\begin{lemma}\label{lem:chow-descent}
Let $k'/k$ be a finite extension of finite fields. Suppose that
$\mathsf h(X_{k'})$ is a summand of a finite direct sum of motives
$\mathsf h(A_j)\otimes\mathbb L^{r_j}$, where the $A_j$ are abelian
varieties over $k'$ and $r_j\ge0$. Then $\mathsf h(X)$ is a summand
of a finite direct sum of abelian, Artin, and Lefschetz motives over
$k$. This holds for Chow retractions and also for correspondences
whose retraction identities hold in every $\ell$-adic realization.
\end{lemma}

\begin{proof}
Regard a $k'$-variety as a $k$-variety and a correspondence over $k'$
as a correspondence over $k$, by extension by zero along the open
and closed inclusion $Y\times_{k'}Z\subset Y\times_k Z$.
This defines induction of motives and preserves composition: after
base change to $k'$, it is the direct sum of the conjugate
correspondences, with zero off the diagonal.

For $A/k'$, put $B=\operatorname{Res}_{k'/k}A$. This is an abelian
variety: the Weil restriction exists for quasiprojective schemes
\citep[\S4.6, Proposition 4.6.3, pp.~110--111]{Poonen2017}, and after
the separable extension $k'/k$ it is the product of the conjugates
of $A$. Projection onto the identity factor and insertion using the
origins in the other factors split $\mathsf h(A)$ from
$\mathsf h(B_{k'})$. Induction therefore splits
$\operatorname{Ind}_{k'/k}\mathsf h(A)$ from
$\mathsf h(B)\otimes\mathsf h(\operatorname{Spec}k')$.
Finally the finite map $X_{k'}\to X$ gives a retraction after dividing
its pushforward by $[k':k]$. Composing these splittings proves the
claim. All identities used are correspondence identities, so the
same argument applies to realization-level retractions.
\end{proof}

\subsection{Kummer surfaces and separated cubic fourfolds}

The following two families admit the required correspondences
directly from their geometry.

\begin{proposition}\label{prop:kummer}
Let $A$ be an abelian surface over a finite field of characteristic
different from two, and let $K=\operatorname{Km}(A)$ be its Kummer
surface. There is a Chow retraction
\[
 \mathsf h(K)\ \mid\
 \mathsf h(A)\oplus\mathsf h(A[2])\otimes\mathbb L.
\]
If $A=E_1\times E_2$, all powers of $K$ satisfy the Tate conjecture
in every codimension over every finite extension. Its formal Brauer
height is one if both elliptic curves are ordinary, and two if exactly
one is ordinary and the other supersingular.
\end{proposition}

\begin{proof}
Blow up the finite etale subscheme $A[2]$. The quotient of
$\widetilde A=\operatorname{Bl}_{A[2]}A$ by the lifted involution
is the smooth Kummer surface, and the quotient map has degree two.
This is the usual Kummer construction
\citep[Chapter 1, Example 1.3(iii), author draft pp.~8--9]{Huybrechts2016}.
Lemma~\ref{lem:retractions} gives the asserted Chow retraction.

For the product case, every power of the right-hand motive is a
summand of sums of motives of products of elliptic curves, with
Artin and Lefschetz factors. Tate classes on products of elliptic
curves are generated by divisor classes; see
\citet[Appendix A.7, first example, pp.~158--159]{Milne2001}.
An Artin factor becomes a sum of units after a finite extension, and
averaging cycles over the finite Galois group descends the Tate
statement. Taking tensor powers of the retraction proves it for $K$.

For the height calculation, the degree-two pullback identifies
$H^2(K)$ with $H^2(A)\oplus H^0(A[2])(-1)$; it is injective and
both sides have dimension $22$. The same holds for crystalline
cohomology. The Artin summand contributes sixteen slopes equal to
one, whether or not the two-torsion points are rational. The slopes
of an ordinary elliptic curve are $0,1$,
and those of a supersingular elliptic curve are $1/2,1/2$.
The Kunneth formula therefore gives
\[
 \begin{array}{c|c}
 (E_1,E_2)&\text{slopes of }H^2(K)\\ \hline
 \text{ordinary, ordinary}&0,\ 1^{20},\ 2\\
 \text{ordinary, supersingular}&(1/2)^2,\ 1^{18},\ (3/2)^2.
 \end{array}
\]
Here a superscript denotes multiplicity. The finite-height K3
Newton polygon has slopes $1-1/h$ and $1+1/h$, each with
multiplicity $h$, and $1$ with multiplicity $22-2h$;
see \citet[Chapter 18, \S3.2--3.3, author draft pp.~399--401]{Huybrechts2016}.
Thus the displayed heights are respectively one and two.
\end{proof}

\begin{proposition}\label{prop:cubic-join}
Let $k$ be a finite field of characteristic different from three.
Let $f$ and $g$ be ternary cubic forms defining smooth plane cubic
curves $E_f,E_g$, and set
\[
 X:\quad f(x_0,x_1,x_2)+g(y_0,y_1,y_2)=0\ \subset\mathbf P^5.
\]
Then $X$ is smooth. After a finite extension, its Chow motive is a
summand of
\begin{equation}\label{eq:cubic-target}
 (\mathbf1\oplus\mathbb L^{\oplus7}\oplus\mathbb L^2)^{\otimes2}
       \ \oplus\ \mathsf h(E_f\times E_g)\otimes\mathbb L.
\end{equation}
Over $k$, it is a summand of abelian, Artin, and Lefschetz motives.
Every power of $X$ satisfies the Tate conjecture in every codimension
over every finite extension.
\end{proposition}

\begin{proof}
Euler's identity and $3\ne0$ show that the gradient of $f$ has
no zero at a nonzero vector, and likewise for $g$. This proves
smoothness of $X$ and of the cubic surfaces
\[
 S_f:\ f(x)+u^3=0,\qquad S_g:\ g(y)+v^3=0.
\]
On their product consider
\[
 ((x:u),(y:v))\dashrightarrow(xv:-uy).
\]
Its image lies in $X$, since
$f(xv)+g(-uy)=v^3(-u^3)-u^3(-v^3)=0$.
The base locus is $E_f\times E_g$, and locally its base ideal is
$(u,v)$: some coordinate of $x$ and some coordinate of $y$ are units.
Consequently its blowup
\[
 V=\operatorname{Bl}_{E_f\times E_g}(S_f\times S_g)
\]
resolves the map. Above a point $(a:b)\in X$ with $f(a)\ne0$,
the preimages are
\[
 ((a:u),(-b/u:1)),\qquad u^3=-f(a).
\]
There are three geometric preimages and the equation is separable.
Thus $V\to X$ is generically finite of degree three. This is the
degree-three instance of the construction of
\citet[Remark 1.10, p.~102]{ShiodaKatsura1979}.

Over a finite extension each cubic surface is a plane blown up at
six points, giving
$\mathsf h(S_f)\simeq\mathsf h(S_g)\simeq
\mathbf1\oplus\mathbb L^{\oplus7}\oplus\mathbb L^2$.
The classification used here is
\citet[Theorem 9.4.4, p.~276]{Poonen2017}.
Lemma~\ref{lem:retractions} now gives \eqref{eq:cubic-target},
with retraction $\frac13\phi_*$. Over a further finite extension
the plane cubics have origins and are elliptic curves.
Lemma~\ref{lem:chow-descent} descends the retraction to $k$.
The Tate assertion follows from the elliptic-product theorem exactly
as in Proposition~\ref{prop:kummer}.
\end{proof}

These are explicit constructions for the stated families. In
particular Proposition~\ref{prop:cubic-join} does not cover an
arbitrary smooth cubic fourfold.

\subsection{The exact mixed correspondence condition}

Fix a finite field $k=\F_q$. A \emph{target with the full Tate
property} will mean a motive in the tensor category generated by
a specified finite collection of abelian varieties, finite etale
schemes, and Lefschetz motives, such that every product of the
abelian generators satisfies the Tate conjecture in every
codimension after every finite extension. This hypothesis is not
asserted for arbitrary collections of abelian varieties.
In such a category, cycle classes span all Frobenius-equivariant
maps between realizations: apply Kunneth and Poincare duality to
the product presenting the two motives, and then apply their
algebraic projectors. Duals and Tate twists give the same statement
for maps to $\mathbb L^a$.

\begin{theorem}\label{thm:mixed-condition}
Let $X/k$ be a finite-height K3 surface or a smooth cubic fourfold,
and let $d=\dim X$. Suppose that a target with the full Tate property
has a finite presentation
\[
 M=(T,e),\qquad T=\bigoplus_{\nu=1}^s\mathsf h(Y_\nu),
\]
where $e$ is a rational algebraic projector in all realizations,
the $Y_\nu$ are products of the target generators and projective
spaces, and $M$ has the same Frobenius characteristic polynomial
as $X$ in every cohomological degree. The following are equivalent.
\begin{enumerate}
\item There are rational algebraic correspondences
$f:\mathsf h(X)\to M$ and $g:M\to\mathsf h(X)$ that are inverse
in every $\ell$-adic realization, $\ell\ne p$.
\item For one prime $\ell_0\ne p$, the projected map
\begin{equation}\label{eq:mixed-map}
 \bigoplus_\nu\operatorname{CH}^{d}(X\times Y_\nu)_{\Q}
       \otimes\Q_{\ell_0}
 \longrightarrow
 \operatorname{Hom}_{F}^{\mathrm{gr}}
       (H^*_{\ell_0}(X),H^*_{\ell_0}(M)),\quad z\longmapsto e_*z_*,
\end{equation}
is surjective.
\item For one prime $\ell_0\ne p$, the Tate conjecture in
codimension $d$ holds for each of the finitely many products
$X\times Y_\nu$.
\end{enumerate}
Here the maps on the right of \eqref{eq:mixed-map} preserve degree
and commute with Frobenius. If these conditions hold at one prime,
they hold at every prime different from $p$.
\end{theorem}

\begin{proof}
Frobenius is semisimple on the cohomology in question. For K3
surfaces, see the semisimplicity input in
\citet[Lemma 10.4 and Remark 10.5, p.~54]{IIK2021}; for cubic
fourfolds, see \citet[introduction, author preprint p.~1]{Zarhin2004}.
Abelian varieties and the other target generators also have
semisimple Frobenius. Thus equality of the degreewise characteristic
polynomials supplies an isomorphism at $\ell_0$.

Assume (2). Lift such an isomorphism to an algebraic correspondence
$f_0$ with $\Q_{\ell_0}$ coefficients. Its inverse is algebraic
with these coefficients as well, for the following reason, which
does not presume any mixed Tate statement in the reverse direction.
Transport the Poincare duality form of $X$ by $f_0$ to a
Frobenius-equivariant pairing
$H^*(M)\otimes H^*(M)\to\Q_{\ell_0}(-d)$.
The full Tate property of the target lifts this pairing to a
correspondence. Compose it with the dual of $f_0$ and the inverse
of the canonical motivic duality of $X$. This gives a correspondence
$g_0:M\to\mathsf h(X)$ realizing $f_0^{-1}$.

Write $f_0,g_0$ as finite $\Q_{\ell_0}$-linear combinations of
rational cycles. Nonvanishing of the determinant of the composition
is an open condition on their coefficient tuples. Since $\Q$ is
dense in $\Q_{\ell_0}$, there are rational combinations $f,g$ for
which $u=gf$ is invertible at $\ell_0$.

Both types of $X$ have zero odd cohomology. For each $r\ge1$,
the correspondence trace formula therefore reads
\begin{equation}\label{eq:trace-certificate}
 \operatorname{Tr}(u_*^r\mid H^*_\ell(X))
           =\deg(\Delta_X\cdot u^r)\in\Q.
\end{equation}
The right side is independent of $\ell$. The formula follows
directly by expanding the diagonal in dual Kunneth bases; see
\citet[Lemmas 45.9.7 and 45.10.3, Tags 0FHH and 0FHR]{Stacks}.
Newton identities give a common rational characteristic polynomial
$P_u(T)$ for $u_*$, of degree $\dim H^*(X)$. Its constant term
is nonzero, because it is nonzero at $\ell_0$. Cayley--Hamilton
expresses $u_*^{-1}$ as a rational polynomial $R(u_*)$ at every
prime. Replace $g$ by $R(u)g$. Then $gf=1$ in all realizations;
equal dimensions give $fg=1$ there as well. This proves (1).

If (1) holds, any Frobenius-equivariant graded map from $H^*(X)$
to $H^*(Y_\nu)$ is the composite of $f$ with a map from $M$ to
$\mathsf h(Y_\nu)$. The full Tate property lifts the latter map.
Kunneth and Poincare duality identify all these maps with
codimension-$d$ Tate classes on $X\times Y_\nu$. Thus (1) implies
(3) at every prime. Finally (3) implies (2) by projection with $e$.
\end{proof}

The condition can be restricted to the remaining summand.
For a cubic fourfold, let $H$ be its hyperplane class and put
\[
 p_{2i}=\tfrac13 H^{4-i}\times H^i\quad(0\le i\le4),
 \qquad p_{\mathrm{prim}}=\Delta_X-\sum_{i=0}^4p_{2i}.
\]
These are mutually orthogonal Chow projectors, since $\deg H^4=3$;
the last one realizes the $22$-dimensional primitive $H^4$.
For a K3 surface, pass to a finite extension defining a basis
$D_1,\ldots,D_\rho$ of the Neron--Severi group over $\Q$.
Choose a rational zero-cycle $o$ of degree one, let
$B=(D_i\cdot D_j)$, and put
\[
 p_0=o\times X,\quad p_4=X\times o,\quad
 p_{\mathrm{NS}}=\sum_{i,j}(B^{-1})_{ij}D_i\times D_j,\quad
 p_{\mathrm{tr}}=\Delta_X-p_0-p_4-p_{\mathrm{NS}}.
\]
The intersection form is nondegenerate, and multiplication of
product correspondences verifies all projector identities in Chow
groups. The complement is
$\mathbf1\oplus\mathbb L^{\oplus\rho}\oplus\mathbb L^2$.

If a target $N=(T,e)$ with the full Tate property has the same
realization as $p_{\mathrm{prim}}$ or $p_{\mathrm{tr}}$, respectively,
the exact remaining condition is the surjectivity of
\[
 z\longmapsto e_*z_*p_*:
 \bigoplus_\nu\operatorname{CH}^{d}(X\times Y_\nu)_{\Q}
            \otimes\Q_{\ell_0}
 \longrightarrow
 \operatorname{Hom}_F(H^*_{\ell_0}(X,p),H^*_{\ell_0}(N)).
\]
Indeed the known complement already has algebraic comparison maps.
Adjoining these maps reduces the assertion to
Theorem~\ref{thm:mixed-condition}. Conversely an isomorphism of the
full motives transports $p$ to an algebraic projector on the target;
the full Tate property supplies the comparison with the specified
$N$. This also explains why the choice of a presentation is harmless.
Here target existence is a hypothesis; no general CM realization
theorem is being assumed without its hypotheses.

\subsection{A target for a weight-two slope pattern}

The ordinary case permits a target construction using only
Honda--Tate theory and a published all-powers Tate theorem. We give
the arithmetic details, including the possible noncommutative
endomorphism algebra.

\begin{proposition}\label{prop:ordinary-packet}
Let $f(T)\in\Z[T]$ be a nonconstant monic polynomial whose roots
are algebraic integers of complex absolute value $q$, where
$q=p^a$. With valuations normalized by $v(q)=1$, suppose that its
roots have slopes $0$ and $2$ each once and slope $1$ otherwise.
Assume that no root is $q$ times a root of unity. Then:
\begin{enumerate}
\item $f$ is irreducible and separable. For a root $\alpha$,
$K=\Q(\alpha)$ is a CM field and $\Q(\alpha^r)=K$ for every $r>0$.
\item There is a geometrically simple abelian variety $B/\F_{q^2}$
with Frobenius $\alpha$. Its endomorphism division algebra $D$ has
center $K$ and degree $d\in\{1,2\}$ over $K$.
\item Every power of $B$ satisfies the Tate conjecture in every
codimension after every finite extension. There is a rational
algebraic homological summand $N$ of $\mathsf h^1(B)^{\otimes2}$
whose Frobenius roots, each once, are
\[
                  \{\sigma(\alpha)^2:\sigma:K\hookrightarrow\overline\Q\}.
\]
All the projectors defining $N$ are the same rational
correspondences for every $\ell\ne p$.
\end{enumerate}
\end{proposition}

\begin{proof}
Choose a root of slope zero. Its irreducible factor also contains
its complex conjugate $q^2/\alpha$, of slope two. That factor has
multiplicity one. Any other factor has all its roots of slope one
at all primes above $p$. For one of its roots $\beta$, the number
$\beta/q$ is an algebraic integer: this is clear away from $p$,
and follows from the slope assertion above $p$. Every complex
conjugate has absolute value one, so Kronecker's theorem makes
$\beta/q$ a root of unity, contrary to the assumption. Thus $f$
is irreducible. Conjugation acts on $K$ by
$\alpha\mapsto q^2/\alpha$, commutes with every embedding, and
has a totally real fixed field. No embedding of $K$ is real,
since a real conjugate would be $q$ or $-q$. Hence $K$ is CM.

For $L=\Q(\alpha^r)$, use the slope-zero embedding of $K$ into
$\overline\Q_p$. Every $L$-embedding of $K$ sends $\alpha$ to a
root with the same $r$th power, hence to another slope-zero root.
There is only one. Separability gives $[K:L]=1$.
The prime of $K$ where $\alpha$ has slope zero has local degree
one: the number of embeddings inducing that prime equals its local
degree and each would give a slope-zero root. Its conjugate also
has local degree one.

Now $\alpha$ is a Weil $q^2$-number. Its slopes relative to $q^2$
are $0,1/2,1$. The Honda--Tate local invariants have denominators
at most two, so the division-algebra degree is $d=1$ or $2$.
The equality $\Q(\alpha^r)=K$ and constancy of these invariants
after every extension prove geometric simplicity and show that
all geometric endomorphisms are already present.

For completeness, the all-powers Tate statement can be seen from
the same unique slope-zero prime. A conjugate of $\alpha/q$ has
valuations supported at just the slope-zero and slope-two
embeddings, with values $-1,+1$. If
\[
        \prod_\sigma (\sigma(\alpha)/q)^{m_\sigma}
\]
is a root of unity, applying all conjugate valuations forces
$m_\sigma=m_{\bar\sigma}$ for each conjugate pair. Thus every
Tate eigenmonomial in the exterior algebra on any number of copies
of $H^1(B)$ is a product of degree-two Tate eigenmonomials.
The divisor Tate theorem makes those factors algebraic. Scalar
extension to a splitting field and descent of the span give the
assertion over $\Q_\ell$. This is precisely the second example of
\citet[Appendix A.7, p.~159]{Milne2001}, which also covers $d=2$;
it is not restricted to the commutative-endomorphism case.

It remains to remove the multiplicity $d$ in $H^1(B)$.
An algebraic projector onto $\mathsf h^1(B)$ is obtained by
polynomial interpolation in $[m]^*$, whose eigenvalues on
$H^i(B)$ are $m^i$, for any integer $m>1$.
The separable algebra $K\otimes_\Q K$ has an idempotent
$e_\Delta$ for its multiplication factor $K$; over a splitting
field it keeps exactly the pairs of equal embeddings.
Inside the algebra of endomorphisms of
$\mathsf h^1(B)^{\otimes2}$, its block is
\[
 e_\Delta(D\otimes_\Q D)\simeq D\otimes_K D.
\]
Every local invariant of $D$ has order at most two. The last
algebra therefore has trivial Brauer class and is
$\operatorname{Mat}_{d^2}(K)$. Choose a primitive idempotent $j$
in it. On the equal-embedding block the realization has dimension
$d^2$ at each embedding of $K$, and this matrix algebra acts by
its standard module. The image of $j$ consequently has dimension
one at each embedding and Frobenius $\sigma(\alpha)^2$ there.
The corresponding rational algebraic projector defines $N$.
An opposite-algebra convention for $\mathsf h^1$ changes none of
this argument, since the Brauer class has order at most two.
The asserted idempotences hold modulo homological equivalence;
no lifting to Chow idempotents is asserted here.
\end{proof}

\subsection{Ordinary K3 surfaces and ordinary cubic fourfolds}

We apply the slope-pattern construction after removing the
algebraic part of the middle cohomology.

\begin{proposition}\label{prop:ordinary-target}
Let $X/\F_q$ be an ordinary K3 surface or an ordinary smooth cubic
fourfold, in any characteristic. After a finite extension, there is
an explicitly specified homological motive $M$ with the full Tate
property whose graded Frobenius realization equals that of $X$.
It is built from powers of a single abelian variety, Lefschetz
motives, and direct summands. For this target the remaining
assertion that $\mathsf h(X)\simeq M$ by rational algebraic
correspondences is exactly the mixed condition of
Theorem~\ref{thm:mixed-condition}.
\end{proposition}

\begin{proof}
The divisor Tate theorem for K3 surfaces holds also in
characteristics two and three; it follows, in particular, from
\citet[Theorem 1.5, p.~3, and Lemma 10.3, pp.~53--54]{IIK2021}.
For cubic fourfolds, codimension-two Tate in odd characteristic is
\citet[Theorem 5.14(3), author preprint p.~27]{MadapusiPera2015}.
In characteristic two it is
\citet[Theorem 4.10, author-hosted journal text p.~20]{AuelEtAl2025}.
Together with the semisimplicity references in
Theorem~\ref{thm:mixed-condition}, these statements identify the
entire root-of-unity part of the middle cohomology with algebraic
classes after a finite extension. Enlarge the field so all those
classes are defined, and rename the field $\F_q$.

The middle Frobenius polynomial is rational. Its roots of the
form $q^r\zeta$, with $\zeta$ a root of unity and $r=1$ or $2$,
form a Galois-stable multiset. Removing their factor therefore
leaves a monic integral polynomial. Semisimplicity and Poincare
duality make the pairing on the algebraic eigenspace nondegenerate.
Choose rational algebraic cycles $C_1,\ldots,C_t$ whose classes
form a basis of this eigenspace at one prime, and put
$G=(C_i\cdot C_j)$. The rational matrix $G$ is nonsingular, so
the classes remain independent at every prime; the characteristic
polynomial shows that they span the same-dimensional eigenspace
there. Their Chow projector is
$\sum_{i,j}(G^{-1})_{ij}C_i\times C_j$.
Its complement realizes exactly the remaining polynomial.

For an ordinary K3 surface the slopes of $H^2$ are
$0,1^{20},2$. Removing the algebraic factor removes only slope-one
roots. The remaining polynomial meets
Proposition~\ref{prop:ordinary-packet}. After the quadratic
extension $\F_{q^2}$, its nonalgebraic realization is therefore
that of the motive $N$ constructed there. Adjoin
$\mathbf1\oplus\mathbb L^{\oplus\rho}\oplus\mathbb L^2$.

For an ordinary cubic fourfold the middle slopes are
$1,2^{21},3$; see \citet[\S4, discussion after Remark 4.1,
author preprint p.~6]{Zarhin2004}.
Write a remaining root as $\gamma$ and put $\alpha=\gamma/q$.
This is an algebraic integer: outside $p$, division by $q$ is
harmless; at every prime above $p$, all conjugates of $\gamma$
have valuation at least $v(q)$. The polynomial of these
$\alpha$'s is consequently monic integral, of weight two, with
slopes $0,1,2$ as in Proposition~\ref{prop:ordinary-packet}, and
without roots $q\zeta$. Over $\F_{q^2}$, the motive
$N\otimes\mathbb L$ has roots $q^2\alpha^2=\gamma^2$.
Adjoin the algebraic middle summand and the known outer-degree
Lefschetz summands to obtain $M$.

All the required Lefschetz motives occur in powers of the same
$B$: an ample divisor and a suitably normalized complementary
power split $\mathbb L$ from $\mathsf h(B)$. Hence $M$ has a
finite untwisted presentation using powers of $B$, and
Proposition~\ref{prop:ordinary-packet} supplies its full Tate
property. The exact remaining condition is now
Theorem~\ref{thm:mixed-condition}. If the comparison
correspondences are constructed over an extension, induction,
Weil restriction, and trace as in Lemma~\ref{lem:chow-descent}
give a comparison with an abelian--Artin target over the original
field. The conjugates of $B$ have the same Frobenius polynomial
and are isogenous, so this descent does not introduce an
unproved mixed Tate assertion among unrelated abelian varieties.
\end{proof}

The target construction and the mixed comparison are different
conclusions. The proposition proves the former. It does not
construct the missing mixed cycles for an arbitrary ordinary K3
surface or ordinary cubic fourfold. The known Tate theorem for
$X\times X$ in the K3 case
\citep[Theorem 1.5]{IIK2021}, and for all self-powers in the
ordinary cubic case \citep[Corollary 5.2, author preprint p.~8]{Zarhin2004},
do not by themselves imply Tate for $X\times Y_\nu$ with the
specified abelian targets. The retractions of
Propositions~\ref{prop:kummer} and \ref{prop:cubic-join} supply
such cycles for their concrete families; the general mixed
condition remains explicit.
